%%%%%%
%
% $Author: Sadegh Naderi $
% $Datum: 2023-10-15  $
% $Pfad: ML23-01-Keyword-Spotting-with-an-Arduino-Nano-33-BLE-Sense\Presentations\SadeghNaderi\LiteratureReview\slides\ContentLiteratureReview.tex $
% $Version: 3.0 $
% $Review Date: 2023-12-16 $
%
%
% !TeX encoding = utf8
% !TeX root = Rename
%
%%%%%%


\Mysection{Praxiswissen in der Messtechnik}

\STANDARD{Beschreibung}
{ 
	Das Buch behandelt grundlegende und praxisnahe Themen der Messtechnik. Es erklärt, wie Messgrößen erfasst, bewertet und technisch verarbeitet werden. Die Quelle eignet sich als Grundlage für Messfehler, Genauigkeit, Messverfahren und die Auswahl geeigneter messtechnischer Vorgehensweisen. Die Quelle wird mit \cite{Helbig:2021} belegt.
	
	\bigskip
	
	\textbf{Limitation:} Die Quelle ist allgemein auf Messtechnik ausgelegt und beschreibt keine einzelnen Bauteile oder konkreten Schaltungsaufbauten im Detail.
	
	\scriptsize{\textbf{Keywords:} Messtechnik, Messabweichung, Messgenauigkeit, Messverfahren, Sensorauswertung}
}


\Mysection{Messelektronik und Sensoren}

\STANDARD{Beschreibung}
{ 
	Das Buch beschreibt Grundlagen der Messelektronik, Sensoren sowie analoge und digitale Signalverarbeitung. Es hilft dabei zu verstehen, wie elektrische Messgrößen, Sensorsignale und Schaltsignale aufgenommen, aufbereitet und ausgewertet werden. Die Quelle wird mit \cite{Bernstein:2024} belegt.
	
	\bigskip
	
	\textbf{Limitation:} Die Quelle ist sehr umfangreich und behandelt viele Sensorarten allgemein. Konkrete Hinweise zu einzelnen Mikrocontroller-Boards stehen nicht im Vordergrund.
	
	\scriptsize{\textbf{Keywords:} Messelektronik, Sensorsignal, Signalaufbereitung, digitale Signalverarbeitung, Messkette}
}


\Mysection{Elektrotechnik und Elektronik für Maschinenbauer}

\STANDARD{Beschreibung}
{ 
	Das Buch vermittelt Grundlagen der Elektrotechnik und Elektronik aus Sicht des Maschinenbaus. Es eignet sich als Basisquelle für Spannung, Strom, Widerstand, elektrische Schaltungen und elektronische Bauteile. Die Quelle wird mit \cite{Hering:2012} belegt.
	
	\bigskip
	
	\textbf{Limitation:} Das Buch liefert Grundlagenwissen, aber keine detaillierte Schaltungsauslegung für einzelne elektronische Baugruppen.
	
	\scriptsize{\textbf{Keywords:} elektrische Grundlagen, Gleichspannung, Stromfluss, Widerstand, Grundschaltung}
}


\Mysection{Mikrosysteme}

\STANDARD{Beschreibung}
{ 
	Das Buch beschreibt Mikrosysteme und deren technische Grundlagen. Es zeigt, wie kleine elektronische und mechanische Systeme miteinander kombiniert werden können und unterstützt das Verständnis kompakter mechatronischer Baugruppen. Die Quelle wird mit \cite{zielke:2025} belegt.
	
	\bigskip
	
	\textbf{Limitation:} Die Quelle ist grundlagenorientiert und beschreibt keine konkreten Arduino- oder CNC-Aufbauten.
	
	\scriptsize{\textbf{Keywords:} Mikrosystemtechnik, mechatronische Baugruppe, Sensor-Aktor-System, Miniaturisierung, Systemintegration}
}


\Mysection{Sensortechnik}

\STANDARD{Beschreibung}
{ 
	Das Buch behandelt Grundlagen und Anwendungen der Sensortechnik. Es erklärt den Zusammenhang zwischen physikalischer Größe, elektrischem Signal und technischer Auswertung. Dadurch eignet es sich für die allgemeine Einordnung von Sensoren und Schaltern. Die Quelle wird mit \cite{traenkler:2014} belegt.
	
	\bigskip
	
	\textbf{Limitation:} Die Quelle betrachtet Sensortechnik allgemein und beschreibt kein einzelnes Endschaltermodell.
	
	\scriptsize{\textbf{Keywords:} Sensortechnik, Positionsrückmeldung, Zustandserfassung, Sensorsignal, Messprinzip}
}


\Mysection{Creality Endschalter von 3DJAKE}

\STANDARD{Beschreibung}
{ 
	Die Produktquelle beschreibt einen mechanischen Endschalter, wie er in 3D-Druckern und ähnlichen Bewegungssystemen eingesetzt werden kann. Sie ist geeignet, um Bauform, Einsatzgebiet und grundsätzliche Funktion eines solchen Schalters einzuordnen. Die Quelle wird mit \cite{3DJAKE:2026} belegt.
	
	\bigskip
	
	\textbf{Limitation:} Es handelt sich um eine Produktseite. Sie liefert nur begrenzt technische Hintergrundinformationen und ersetzt kein vollständiges Datenblatt.
	
	\scriptsize{\textbf{Keywords:} Endschalter, Creality, mechanischer Schalter, Positionsbegrenzung, Referenzsignal}
}


\Mysection{Drucktaster von Eaton: Ein Leitfaden}

\STANDARD{Beschreibung}
{ 
	Die Website erklärt allgemein, was ein Drucktaster ist, wie er funktioniert und wofür er eingesetzt wird. Sie beschreibt Drucktaster als Bedien- und Eingabeelemente, mit denen elektrische Schaltvorgänge ausgelöst werden können. Die Quelle wird mit \cite{RS:2026} belegt.
	
	\bigskip
	
	\textbf{Limitation:} Die Quelle ist als Produktratgeber aufgebaut und weniger wissenschaftlich belastbar als ein Fachbuch oder Datenblatt.
	
	\scriptsize{\textbf{Keywords:} Drucktaster, Tastsignal, Bedienelement, Schließerkontakt, Eingabeauswertung}
}


\Mysection{Arcade-Buttons mit Mikroschalter}

\STANDARD{Beschreibung}
{ 
	Die Quelle beschreibt Arcade-Buttons mit Mikroschalter. Sie ist relevant, wenn mechanische Taster als einfache digitale Eingabeelemente betrachtet werden. Solche Taster können für Bedienfelder und Benutzeraktionen eingesetzt werden. Die Quelle wird mit \cite{Joy-it:2024} belegt.
	
	\bigskip
	
	\textbf{Limitation:} Die Quelle enthält nur wenige technische Details und ist produktbezogen. Für genaue elektrische Grenzwerte sind zusätzliche Datenblätter sinnvoll.
	
	\scriptsize{\textbf{Keywords:} Arcade-Button, Mikroschalter, Drucktaster, Bedienfeld, Schaltkontakt}
}


\Mysection{Einführung in die elektronische Schaltungstechnik}

\STANDARD{Beschreibung}
{ 
	Das Buch behandelt Grundlagen elektronischer Schaltungen. Es unterstützt das Verständnis von Bauteilen, Schaltungslogik, Signalpfaden und elektrischer Verbindungstechnik. Die Quelle wird mit \cite{Alles:2024} belegt.
	
	\bigskip
	
	\textbf{Limitation:} Die Quelle ist eine allgemeine Einführung und ersetzt keine konkreten Datenblätter einzelner Bauteile.
	
	\scriptsize{\textbf{Keywords:} Schaltungstechnik, Bauteilauswahl, Signalpfad, elektrische Verbindung, Grundschaltung}
}


\Mysection{Halbleiter-Schaltungstechnik}

\STANDARD{Beschreibung}
{ 
	Das Buch beschreibt Halbleiterbauelemente und elektronische Schaltungen sehr detailliert. Es eignet sich besonders als Hintergrundquelle für MOSFETs, Treiberschaltungen, Schaltvorgänge und Verlustleistung. Die Quelle wird mit \cite{Tietze:2023} belegt.
	
	\bigskip
	
	\textbf{Limitation:} Die Quelle ist fachlich sehr tiefgehend und für einfache technische Beschreibungen teilweise umfangreicher als notwendig.
	
	\scriptsize{\textbf{Keywords:} Halbleiterbauelement, MOSFET-Schaltung, Leistungsschalter, Gate-Ansteuerung, Verlustleistung}
}


\Mysection{Grundkurs Leistungselektronik}

\STANDARD{Beschreibung}
{ 
	Das Buch behandelt Grundlagen der Leistungselektronik. Es erklärt, warum größere elektrische Lasten nicht direkt über Logikausgänge geschaltet werden sollten und welche Rolle Leistungsschalter, Schutzfunktionen und Verlustleistung spielen. Die Quelle wird mit \cite{Specovius:2025} belegt.
	
	\bigskip
	
	\textbf{Limitation:} Die Quelle beschreibt Leistungselektronik allgemein und keinen exakt vorgegebenen Schaltungsaufbau.
	
	\scriptsize{\textbf{Keywords:} Leistungselektronik, Lastschaltung, Schalten, Verlustleistung, Leistungsstufe}
}


\Mysection{Mechatronische Systeme}

\STANDARD{Beschreibung}
{ 
	Das Buch beschreibt Grundlagen mechatronischer Systeme. Es ist hilfreich für Systeme, bei denen Mechanik, Elektronik, Sensorik, Aktorik und Software zusammenwirken. Die Quelle wird mit \cite{Isermann:2008} belegt.
	
	\bigskip
	
	\textbf{Limitation:} Die Quelle behandelt mechatronische Systeme allgemein und liefert keine Bauanleitung für einzelne Geräte.
	
	\scriptsize{\textbf{Keywords:} mechatronisches System, Sensor-Aktor-Kopplung, Steuerungstechnik, Systemstruktur, Aktorik}
}


\Mysection{Leistungselektronik}

\STANDARD{Beschreibung}
{ 
	Das Handbuch behandelt Leistungselektronik ausführlich. Es eignet sich als Hintergrundquelle für elektrische Leistungsstufen, Schaltvorgänge und die Ansteuerung leistungsstärkerer Verbraucher. Die Quelle wird mit \cite{Zach:2022} belegt.
	
	\bigskip
	
	\textbf{Limitation:} Die Quelle ist breit angelegt und beschreibt keine einfache, direkt nachbaubare Mikrocontroller-Schaltung.
	
	\scriptsize{\textbf{Keywords:} Leistungselektronik, Leistungsstufe, Lastschaltung, Stromversorgung, Schaltbetrieb}
}


\Mysection{Schaltnetzteile und ihre Peripherie}

\STANDARD{Beschreibung}
{ 
	Das Buch behandelt Schaltnetzteile, Spannungsregelung und deren Peripherie. Es ist relevant, wenn Baugruppen mit unterschiedlichen Versorgungsspannungen betrachtet werden. Die Quelle wird mit \cite{Schlienz:2020} belegt.
	
	\bigskip
	
	\textbf{Limitation:} Die Quelle erklärt Netzteile allgemein und beschreibt nicht direkt ein konkretes Netzteilmodell.
	
	\scriptsize{\textbf{Keywords:} Schaltnetzteil, Spannungsregelung, DC-DC-Wandler, Versorgungsstabilität, EMV}
}


\Mysection{A4988 Microstepping Driver}

\STANDARD{Beschreibung}
{ 
	Das Datenblatt beschreibt den A4988-Schrittmotortreiber. Es liefert Informationen zu Microstepping, Strombegrenzung, Schutzfunktionen und den benötigten Steuersignalen. Die Quelle wird mit \cite{Allegro:A4988} belegt.
	
	\bigskip
	
	\textbf{Limitation:} Das Datenblatt ist technisch und setzt Grundwissen voraus. Die praktische Einstellung und thermische Auslegung müssen zusätzlich betrachtet werden.
	
	\scriptsize{\textbf{Keywords:} A4988, Schrittmotortreiber, Microstepping, Strombegrenzung, Steuersignal}
}


\Mysection{MOSFET PWM-Schaltmodul}

\STANDARD{Beschreibung}
{ 
	Das Datenblatt beschreibt ein MOSFET-PWM-Schaltmodul zum Schalten oder Regeln elektrischer Lasten über ein Steuersignal. Es eignet sich zur allgemeinen Beschreibung einer Schnittstelle zwischen Logiksignal und Last. Die Quelle wird mit \cite{JoyITMOSFET:2025} belegt.
	
	\bigskip
	
	\textbf{Limitation:} Die tatsächliche Belastbarkeit hängt von Versorgungsspannung, Strom, Kühlung und Verdrahtung ab.
	
	\scriptsize{\textbf{Keywords:} MOSFET, PWM, Schaltmodul, Laststeuerung, Leistungselektronik}
}


\Mysection{IRF9540N P-Channel Power MOSFET}

\STANDARD{Beschreibung}
{ 
	Das Datenblatt beschreibt den P-Kanal-Leistungs-MOSFET IRF9540N. Es liefert wichtige Grenzwerte zu Spannung, Strom, Verlustleistung und thermischer Belastung. Die Quelle wird mit \cite{IRF9540N:2025} belegt.
	
	\bigskip
	
	\textbf{Limitation:} Ein Datenblatt zeigt Grenzwerte des Bauteils. Die Eignung hängt zusätzlich von Schaltung, Ansteuerung und Wärmeabfuhr ab.
	
	\scriptsize{\textbf{Keywords:} IRF9540N, P-Kanal-MOSFET, Leistungs-MOSFET, Lastschaltung, Datenblatt}
}


\Mysection{Creality Endschalter von bastelgarage}

\STANDARD{Beschreibung}
{ 
	Die Produktquelle beschreibt einen Creality-Endschalter. Sie ist geeignet, um den Einsatz mechanischer Endschalter zur Erkennung von Positionen oder Grenzlagen allgemein zu beschreiben. Die Quelle wird mit \cite{bastelgarage:2026} belegt.
	
	\bigskip
	
	\textbf{Limitation:} Es handelt sich um eine Produktseite mit begrenzten technischen Detailangaben. Für elektrische Kennwerte ist ein Datenblatt geeigneter.
	
	\scriptsize{\textbf{Keywords:} Endschalter, Creality, Grenzlagenerkennung, Positionssignal, Mikroschalter}
}


\Mysection{setup() - Arduino Reference}

\STANDARD{Beschreibung}
{ 
	Die Arduino-Referenz beschreibt die Funktion setup(). Sie erklärt, dass dieser Programmteil beim Start eines Sketches einmal ausgeführt wird und typischerweise zur Initialisierung verwendet wird. Die Quelle wird mit \cite{arduino_setup} belegt.
	
	\bigskip
	
	\textbf{Limitation:} Die Quelle erklärt nur den einzelnen Befehl und keine vollständige Programmstruktur.
	
	\scriptsize{\textbf{Keywords:} setup, Arduino-Sketch, Initialisierung, Programmstart, Arduino-IDE}
}


\Mysection{loop() - Arduino Reference}

\STANDARD{Beschreibung}
{ 
	Die Arduino-Referenz beschreibt die Funktion loop(). Sie erklärt, dass dieser Programmteil nach setup() fortlaufend wiederholt wird und damit die dauerhafte Programmausführung abbildet. Die Quelle wird mit \cite{arduino_loop} belegt.
	
	\bigskip
	
	\textbf{Limitation:} Die Quelle erklärt nur den einzelnen Befehl und keine konkrete Programmlogik.
	
	\scriptsize{\textbf{Keywords:} loop, Arduino-Sketch, Programmschleife, Dauerbetrieb, Ablaufstruktur}
}


\Mysection{pinMode() - Arduino Reference}

\STANDARD{Beschreibung}
{ 
	Die Arduino-Referenz beschreibt pinMode(). Mit diesem Befehl wird festgelegt, ob ein Pin als Eingang oder Ausgang verwendet wird. Die Quelle wird mit \cite{arduino_pinmode} belegt.
	
	\bigskip
	
	\textbf{Limitation:} Die Quelle beschreibt die Funktion allgemein und nicht die elektrische Belastbarkeit einzelner Pins.
	
	\scriptsize{\textbf{Keywords:} pinMode, Pin-Konfiguration, Eingang, Ausgang, Arduino}
}


\Mysection{digitalWrite() - Arduino Reference}

\STANDARD{Beschreibung}
{ 
	Die Arduino-Referenz beschreibt digitalWrite(). Der Befehl setzt einen digitalen Ausgang auf HIGH oder LOW und kann dadurch digitale Signale ausgeben. Die Quelle wird mit \cite{arduino_digitalwrite} belegt.
	
	\bigskip
	
	\textbf{Limitation:} Die Quelle erklärt den Softwarebefehl, ersetzt aber keine Betrachtung der zulässigen Strombelastung eines Ausgangs.
	
	\scriptsize{\textbf{Keywords:} digitalWrite, digitaler Ausgang, HIGH, LOW, Schaltsignal}
}


\Mysection{Digital Pins und INPUT PULLUP - Arduino Reference}

\STANDARD{Beschreibung}
{ 
	Die Arduino-Referenz beschreibt digitale Pins und den Modus INPUT PULLUP. Dieser Modus nutzt einen internen Pull-up-Widerstand, wodurch einfache Taster- oder Schaltereingänge ohne externen Pull-up-Widerstand möglich sein können. Die Quelle wird mit \cite{arduino_inputpullup} belegt.
	
	\bigskip
	
	\textbf{Limitation:} Die Quelle beschreibt das Prinzip allgemein. Die korrekte Verschaltung und Signalinterpretation müssen zusätzlich beachtet werden.
	
	\scriptsize{\textbf{Keywords:} INPUT PULLUP, Pull-up-Widerstand, digitaler Eingang, Tastsignal, Arduino}
}


\Mysection{digitalRead() - Arduino Reference}

\STANDARD{Beschreibung}
{ 
	Die Arduino-Referenz beschreibt digitalRead(). Der Befehl liest den Zustand eines digitalen Eingangs und gibt HIGH oder LOW zurück. Die Quelle wird mit \cite{arduino_digitalread} belegt.
	
	\bigskip
	
	\textbf{Limitation:} Die Quelle erklärt den Softwarebefehl, aber nicht die elektrische Entprellung oder Störsicherheit von Eingangssignalen.
	
	\scriptsize{\textbf{Keywords:} digitalRead, digitaler Eingang, HIGH, LOW, Signalabfrage}
}


\Mysection{Serial.print() - Arduino Reference}

\STANDARD{Beschreibung}
{ 
	Die Arduino-Referenz beschreibt Serial.print(). Der Befehl wird verwendet, um Werte oder Texte über die serielle Schnittstelle auszugeben. Die Quelle wird mit \cite{arduino_serialprint} belegt.
	
	\bigskip
	
	\textbf{Limitation:} Die Quelle beschreibt die Ausgabe allgemein und keine vollständige Kommunikationsstrategie.
	
	\scriptsize{\textbf{Keywords:} Serial.print, serielle Ausgabe, Debugging, Monitor, Textausgabe}
}


\Mysection{Serial.begin() - Arduino Reference}

\STANDARD{Beschreibung}
{ 
	Die Arduino-Referenz beschreibt Serial.begin(). Mit diesem Befehl wird die serielle Kommunikation gestartet und eine Übertragungsgeschwindigkeit festgelegt. Die Quelle wird mit \cite{arduino_serialbegin} belegt.
	
	\bigskip
	
	\textbf{Limitation:} Die Quelle erklärt die Initialisierung der seriellen Schnittstelle, aber keine weitergehenden Protokolle.
	
	\scriptsize{\textbf{Keywords:} Serial.begin, serielle Kommunikation, Baudrate, Initialisierung, Arduino}
}


\Mysection{millis() - Arduino Reference}

\STANDARD{Beschreibung}
{ 
	Die Arduino-Referenz beschreibt millis(). Der Befehl gibt die seit Programmstart vergangene Zeit in Millisekunden zurück und ermöglicht zeitabhängige Abläufe ohne blockierende Pausen. Die Quelle wird mit \cite{arduino_millis} belegt.
	
	\bigskip
	
	\textbf{Limitation:} Die Quelle beschreibt die Funktion allgemein. Überlaufverhalten und saubere Zeitlogik müssen bei längeren Laufzeiten beachtet werden.
	
	\scriptsize{\textbf{Keywords:} millis, Zeitmessung, nicht blockierender Ablauf, Millisekunden, Arduino}
}


\Mysection{delay() - Arduino Reference}

\STANDARD{Beschreibung}
{ 
	Die Arduino-Referenz beschreibt delay(). Der Befehl hält ein Programm für eine angegebene Zeit an und ist für einfache Zeitverzögerungen geeignet. Die Quelle wird mit \cite{arduino_delay} belegt.
	
	\bigskip
	
	\textbf{Limitation:} Die Quelle beschreibt eine einfache Wartefunktion. Während delay() läuft, ist das Programm nur eingeschränkt reaktionsfähig.
	
	\scriptsize{\textbf{Keywords:} delay, Wartezeit, Zeitverzögerung, blockierender Ablauf, Arduino}
}


\Mysection{Arduino UNO R3}

\STANDARD{Beschreibung}
{ 
	Die offizielle Arduino-Quelle beschreibt das UNO-R3-Board. Sie nennt wichtige Eigenschaften wie Mikrocontroller, Ein- und Ausgänge, Versorgung, Schnittstellen und grundlegende technische Daten. Die Quelle wird mit \cite{arduino:uno-r3} belegt.
	
	\bigskip
	
	\textbf{Limitation:} Die Quelle beschreibt das Board selbst, aber keine vollständige Verschaltung mit externen Komponenten.
	
	\scriptsize{\textbf{Keywords:} Arduino UNO R3, Mikrocontroller, digitale Pins, analoge Eingänge, Boarddaten}
}


\Mysection{Overview of the Arduino UNO Components}

\STANDARD{Beschreibung}
{ 
	Die Arduino-Quelle erklärt die wichtigsten Bauteile und Anschlüsse des Arduino UNO. Sie hilft bei der Einordnung von Pins, Versorgung, Schnittstellen und zentralen Boardkomponenten. Die Quelle wird mit \cite{arduino:uno-components} belegt.
	
	\bigskip
	
	\textbf{Limitation:} Die Quelle bleibt auf das Board begrenzt und beschreibt keine komplette externe Schaltung.
	
	\scriptsize{\textbf{Keywords:} Arduino-Uno-Anschlüsse, Pinbelegung, Boardkomponenten, Versorgungspins, Schnittstellen}
}


\Mysection{Arduino Language Reference}

\STANDARD{Beschreibung}
{ 
	Die Arduino Language Reference fasst die Sprachelemente und Standardfunktionen der Arduino-Programmierung zusammen. Sie eignet sich als allgemeine Quelle für Programmstruktur, Ein- und Ausgänge, Kommunikation und Zeitfunktionen. Die Quelle wird mit \cite{arduino:language-reference} belegt.
	
	\bigskip
	
	\textbf{Limitation:} Die Referenz erklärt einzelne Befehle und Konzepte, aber keine vollständige Anwendung.
	
	\scriptsize{\textbf{Keywords:} Arduino Reference, Arduino-Sprache, Standardfunktionen, Sketch, Programmierung}
}


\Mysection{Arduino UNO R3 Datasheet}

\STANDARD{Beschreibung}
{ 
	Das Datenblatt enthält technische Kenndaten des Arduino UNO R3. Dazu gehören elektrische Eigenschaften, Schnittstellen, Pinbelegung und Boardinformationen. Die Quelle wird mit \cite{arduino:uno-datasheet} belegt.
	
	\bigskip
	
	\textbf{Limitation:} Das Datenblatt beschreibt das Board, aber nicht automatisch die sichere Auslegung angeschlossener externer Lasten.
	
	\scriptsize{\textbf{Keywords:} Arduino UNO, Datenblatt, Pinbelegung, technische Kenndaten, Mikrocontroller-Board}
}


\Mysection{analogRead() - Arduino Reference}

\STANDARD{Beschreibung}
{ 
	Die Arduino-Referenz beschreibt analogRead(). Der Befehl liest analoge Eingangsspannungen ein und wandelt sie in einen digitalen Zahlenwert um. Die Quelle wird mit \cite{arduino:analogread} belegt.
	
	\bigskip
	
	\textbf{Limitation:} Die Quelle beschreibt die Funktion allgemein. Messgenauigkeit und Signalqualität hängen zusätzlich von Schaltung und Referenzspannung ab.
	
	\scriptsize{\textbf{Keywords:} analogRead, Analogeingang, ADC, Messwert, Arduino}
}


\Mysection{analogWrite() - Arduino Reference}

\STANDARD{Beschreibung}
{ 
	Die Arduino-Referenz beschreibt analogWrite(). Der Befehl erzeugt ein PWM-Signal, das für Helligkeitssteuerung, Drehzahlsteuerung oder Signalansteuerung genutzt werden kann. Die Quelle wird mit \cite{arduino:analogwrite} belegt.
	
	\bigskip
	
	\textbf{Limitation:} Die Quelle beschreibt die Softwarefunktion. Die Eignung für eine Last hängt von der nachgeschalteten Elektronik ab.
	
	\scriptsize{\textbf{Keywords:} analogWrite, PWM, Tastverhältnis, Ausgangssignal, Arduino}
}


\Mysection{Arduino Libraries}

\STANDARD{Beschreibung}
{ 
	Die Arduino-Quelle beschreibt Bibliotheken und deren Nutzung. Bibliotheken erweitern die Arduino-IDE um Funktionen für bestimmte Module, Sensoren, Displays oder Kommunikationsschnittstellen. Die Quelle wird mit \cite{arduinoLibraries:2025} belegt.
	
	\bigskip
	
	\textbf{Limitation:} Die Quelle ist allgemein und beschreibt nicht jede einzelne Bibliothek im Detail.
	
	\scriptsize{\textbf{Keywords:} Arduino-Bibliothek, Softwareerweiterung, Modulsteuerung, IDE, Funktionssammlung}
}


\Mysection{Arduino Cloud}

\STANDARD{Beschreibung}
{ 
	Die Arduino-Quelle beschreibt die Arduino Cloud als Online-Plattform für Geräteverwaltung, Datenanzeige und IoT-Anwendungen. Sie ist für vernetzte oder online verwaltete Anwendungen relevant. Die Quelle wird mit \cite{arduinoCloud:2025} belegt.
	
	\bigskip
	
	\textbf{Limitation:} Die Quelle betrifft Cloud-Funktionen und ist nur dann direkt relevant, wenn Online-Anbindung oder Gerätemanagement betrachtet wird.
	
	\scriptsize{\textbf{Keywords:} Arduino Cloud, IoT, Gerätemanagement, Online-Dashboard, Datenübertragung}
}


\Mysection{MP2307 LM2596 Kis-3r33S DC-DC Converter Buck Module}

\STANDARD{Beschreibung}
{ 
	Das Datenblatt beschreibt ein DC-DC-Abwärtswandlermodul. Solche Module wandeln eine höhere Eingangsspannung in eine niedrigere Ausgangsspannung um und werden zur Spannungsversorgung elektronischer Baugruppen eingesetzt. Die Quelle wird mit \cite{MP2307:Datasheet} belegt.
	
	\bigskip
	
	\textbf{Limitation:} Das Datenblatt enthält technische Daten, aber keine ausführliche didaktische Erklärung zur Dimensionierung einer vollständigen Stromversorgung.
	
	\scriptsize{\textbf{Keywords:} DC-DC-Wandler, Buck Converter, Spannungsversorgung, MP2307, Ausgangsspannung}
}


\Mysection{Arduino Uno Rev3 Datasheet}

\STANDARD{Beschreibung}
{ 
	Die Quelle beschreibt technische Daten des Arduino Uno Rev3. Sie eignet sich zur Absicherung von Angaben zu Boardeigenschaften, Anschlüssen und elektrischen Grunddaten. Die Quelle wird mit \cite{Arduino:Uno} belegt.
	
	\bigskip
	
	\textbf{Limitation:} Die Quelle beschreibt das Mikrocontroller-Board, aber nicht die vollständige Auslegung externer Beschaltungen.
	
	\scriptsize{\textbf{Keywords:} Arduino Uno Rev3, Datenblatt, Mikrocontroller-Board, Pinbelegung, technische Daten}
}


\Mysection{Velleman WPI438 Downloads}

\STANDARD{Beschreibung}
{ 
	Die Downloadseite führt Unterlagen zum Velleman-WPI438-Modul auf. Sie dient als Einstiegspunkt zu technischen Dokumenten und unterstützenden Dateien des Displaymoduls. Die Quelle wird mit \cite{Velleman:WPI438} belegt.
	
	\bigskip
	
	\textbf{Limitation:} Die Seite ist eine Downloadübersicht und liefert allein nur begrenzt technische Erklärung.
	
	\scriptsize{\textbf{Keywords:} Velleman WPI438, OLED, Downloadseite, Displaymodul, technische Unterlagen}
}


\Mysection{Velleman VMA438 Data Sheets}

\STANDARD{Beschreibung}
{ 
	Das Datenblatt beschreibt technische Eigenschaften des VMA438- beziehungsweise WPI438-Displaymoduls. Es ist hilfreich für Anschlüsse, elektrische Daten und grundlegende Modulparameter. Die Quelle wird mit \cite{Velleman:VMA438Datasheet} belegt.
	
	\bigskip
	
	\textbf{Limitation:} Das Datenblatt beschreibt das Modul, aber keine vollständige Einbindung in eine bestimmte Softwarearchitektur.
	
	\scriptsize{\textbf{Keywords:} VMA438, OLED-Datenblatt, Displaymodul, Anschlüsse, elektrische Daten}
}


\Mysection{Velleman WPI438 User Manual}

\STANDARD{Beschreibung}
{ 
	Das Benutzerhandbuch beschreibt die Verwendung des WPI438-Displaymoduls. Es unterstützt das Verständnis von Anschluss, grundlegender Nutzung und praktischer Inbetriebnahme. Die Quelle wird mit \cite{Velleman:WPI438UserManual} belegt.
	
	\bigskip
	
	\textbf{Limitation:} Das Handbuch ist produktbezogen und ersetzt keine allgemeine Einführung in Displaytechnik oder Kommunikationsprotokolle.
	
	\scriptsize{\textbf{Keywords:} WPI438, Benutzerhandbuch, OLED, Displayanschluss, Inbetriebnahme}
}


\Mysection{Computergestützte Audio- und Videotechnik}

\STANDARD{Beschreibung}
{ 
	Das Buch behandelt Grundlagen der computergestützten Audio- und Videotechnik. Es kann ergänzend für digitale Ausgabe, Anzeige und Informationsdarstellung genutzt werden. Die Quelle wird mit \cite{Stotz:2019} belegt.
	
	\bigskip
	
	\textbf{Limitation:} Die Quelle passt nur indirekt zu Steuerungs- und Elektronikthemen und ist für Motor- oder Schaltsignaltechnik weniger zentral.
	
	\scriptsize{\textbf{Keywords:} Displaytechnik, digitale Anzeige, Ausgabemedium, Nutzerinformation, visuelle Rückmeldung}
}


\Mysection{Computerunterstützte Produktion}

\STANDARD{Beschreibung}
{ 
	Das Buch gibt eine Einführung in die computerunterstützte Produktion. Es hilft dabei, den Zusammenhang zwischen digitaler Konstruktion, Steuerdaten und automatisierter Fertigung einzuordnen. Die Quelle wird mit \cite{Hehenberger:2020} belegt.
	
	\bigskip
	
	\textbf{Limitation:} Die Quelle behandelt Produktionssysteme allgemein und geht nicht speziell auf einzelne Steuerungen oder Schneidverfahren ein.
	
	\scriptsize{\textbf{Keywords:} CNC-Fertigung, CAM-Daten, automatisierte Bearbeitung, Steuerdaten, computerunterstützte Produktion}
}


\Mysection{Joy-IT ARD-CNC-KIT1 Datenblatt}

\STANDARD{Beschreibung}
{ 
	Das Datenblatt beschreibt das ARD-CNC-KIT1 beziehungsweise ein CNC-Shield und die zugehörigen technischen Eigenschaften. Es ist hilfreich für die Einordnung von Anschlüssen, Motortreibern und CNC-Signalen. Die Quelle wird mit \cite{JoyIT:CNCShield:Datenblatt} belegt.
	
	\bigskip
	
	\textbf{Limitation:} Das Datenblatt ist produktbezogen und behandelt individuelle Anpassungen oder Fehlersuche nur begrenzt.
	
	\scriptsize{\textbf{Keywords:} CNC-Shield, Joy-IT, Arduino-Shield, Motortreiber, Anschlussbelegung}
}


\Mysection{Joy-IT ARD-CNC-KIT1 Anleitung}

\STANDARD{Beschreibung}
{ 
	Die Anleitung beschreibt den grundlegenden Umgang mit dem ARD-CNC-KIT1. Sie unterstützt die praktische Einordnung von Anschluss, Aufbau und Inbetriebnahme eines CNC-Shields. Die Quelle wird mit \cite{JoyIT:CNCShield:Anleitung} belegt.
	
	\bigskip
	
	\textbf{Limitation:} Die Anleitung beschreibt die Produktnutzung, ersetzt aber keine tiefergehende Betrachtung von Steuerungssoftware oder elektrischer Auslegung.
	
	\scriptsize{\textbf{Keywords:} CNC-Shield, Anleitung, Inbetriebnahme, Verdrahtung, Motortreiber}
}


